% GENERATED FILE. Edit canonical lessons, then run npm run generate:books.
% canonical-source-hash: fnv1a64:af4ad14d995214c1
% canonical-lessons: PA-C176-eh, PA-C176-vich, PA-C176-location-frame

\chapter{This, Here, and In a Place}
\label{ch:pa-pointing-place-a1}
\hlchaptermodality{176}

\begin{chapteropening}
\textbf{By the end of this chapter:} \emph{I can point to this or that, ask where a familiar thing is, and answer that it is here, there, or in a familiar place.}

Ask where a familiar book or shop is, then answer with here, there, or a familiar place plus vich and hai.
\end{chapteropening}

\section[\texorpdfstring{\pa{ਇਹ} (eh)}{ਇਹ (eh)}]{\pa{ਇਹ} (eh) — this, near here}

\label{lesson:PA-C176-eh}



\begin{quote}
Put one familiar contrast back in your voice. Say \textbf{\pa{ਇੱਥੇ}} (\emph{itthe}, here), then \textbf{\pa{ਉੱਥੇ}} (\emph{utthe}, there). Ask \textbf{\pa{ਕਿੱਥੇ}?} (\emph{kitthe?}, where?). Now say \textbf{\pa{ਕਿਤਾਬ}} (\emph{kitāb}, book), and the old distant pointer \textbf{\pa{ਉਹ}} (\emph{oh}, that).
\end{quote}

\subsection*{You'll want to know: \pa{ਇਹ}}
\begin{quote}
\textbf{\pa{ਇਹ}} — \emph{eh} — \textbf{this}, near here
\end{quote}

\begin{quote}
\textbf{\pa{ਉਹ}} — \emph{oh} — \textbf{that}, farther away
\end{quote}

You have known \textbf{\pa{ਉਹ}} from its earlier lesson as \emph{he, she, it}. That lesson also told you it points to \textbf{that}. Nothing new is being added to \textbf{\pa{ਉਹ}} now. The new half of the pair is \textbf{\pa{ਇਹ}}: use \emph{eh} for a nearby thing.

The same near/far choice already lives in \textbf{\pa{ਇੱਥੇ} / \pa{ਉੱਥੇ}}, \emph{here / there}. Keep the pairs aligned:

\noindent
\begin{tabularx}{\linewidth}{@{}>{\raggedright\arraybackslash}X>{\raggedright\arraybackslash}X@{}}
\toprule
\textbf{near} & \textbf{farther away} \\
\midrule
\textbf{\pa{ਇਹ}} — this & \textbf{\pa{ਉਹ}} — that \\
\textbf{\pa{ਇੱਥੇ}} — here & \textbf{\pa{ਉੱਥੇ}} — there \\
\bottomrule
\end{tabularx}

Wiktionary traces \textbf{\pa{ਇਹ}} to Sanskrit \emph{eṣa}, “this”.

\begin{grammarlens}[title={What you've built}]
Keep \textbf{\pa{ਹੈ}} at the end, where Punjabi already taught you to put the verb:

\begin{quote}
\textbf{\pa{ਇਹ} \pa{ਕਿਤਾਬ} \pa{ਹੈ।}} — \emph{eh kitāb hai.} — This is a book.
\end{quote}

\begin{quote}
\textbf{\pa{ਉਹ} \pa{ਕਿਤਾਬ} \pa{ਹੈ।}} — \emph{oh kitāb hai.} — That is a book.
\end{quote}

Only the first word changes. The book and the verb stay still.
\end{grammarlens}

\subsection*{Your turn}
\begin{itemize}
\raggedright
  \item \emph{Say it:} for a nearby thing, \textbf{\pa{ਇਹ}}; for a farther thing, \textbf{\pa{ਉਹ}}
  \item \emph{Say it:} \textbf{\pa{ਇਹ} \pa{ਕਿਤਾਬ} \pa{ਹੈ}}
  \item \emph{Switch:} make it distant — \textbf{\pa{ਉਹ} \pa{ਕਿਤਾਬ} \pa{ਹੈ}}
\end{itemize}

\begin{tcolorbox}[breakable,colback=teal!4,colframe=teal!35!black,title={Before you move on}]
Near thing? (\textbf{\pa{ਇਹ}}, \emph{eh}.) Farther thing? (\textbf{\pa{ਉਹ}}, \emph{oh}.) Near place? (\textbf{\pa{ਇੱਥੇ}}, \emph{itthe}.) Farther place? (\textbf{\pa{ਉੱਥੇ}}, \emph{utthe}.)

Source: \href{https://en.wiktionary.org/wiki/\%E0\%A8\%87\%E0\%A8\%B9}{Wiktionary: \pa{ਇਹ}}.
\end{tcolorbox}
\section[\texorpdfstring{\pa{ਵਿੱਚ} (vich)}{ਵਿੱਚ (vich)}]{\pa{ਵਿੱਚ} (vich) — in, at}

\label{lesson:PA-C176-vich}



\begin{quote}
Near and farther: \textbf{\pa{ਇਹ} / \pa{ਉਹ}}. Now bring back two familiar places: \textbf{\pa{ਘਰ}} (\emph{ghar}, home) and \textbf{\pa{ਸਕੂਲ}} (\emph{sakūl}, school). Ask \textbf{\pa{ਕਿੱਥੇ}?} (\emph{kitthe?}, where?).
\end{quote}

\subsection*{You'll want to know: \pa{ਵਿੱਚ}}
\begin{quote}
\textbf{\pa{ਵਿੱਚ}} — \emph{vich} — \textbf{in, at}
\end{quote}

English puts \emph{in} before a place. Punjabi puts \textbf{\pa{ਵਿੱਚ}} after it. That is why it is called a postposition: the place comes first.

\begin{quote}
\textbf{\pa{ਘਰ} \pa{ਵਿੱਚ}} — \emph{ghar vich} — at home, in the house
\end{quote}

\begin{quote}
\textbf{\pa{ਸਕੂਲ} \pa{ਵਿੱਚ}} — \emph{sakūl vich} — at school, in the school
\end{quote}

Do not add a verb yet. This step owns only the place chunk: noun, then \textbf{\pa{ਵਿੱਚ}}.

\begin{cousinweb}
\textbf{\pa{ਵਿੱਚ}} is inherited through Old Punjabi \emph{vici} and Prakrit \emph{vicca} from a form reconstructed as Sanskrit \emph{vīcya}, “middle” (CDIAL 12042). Hindi \emph{bīc}, “middle”, is a cousin. The old idea of being in the middle became Punjabi's ordinary word for being \textbf{in} a place.
\end{cousinweb}

\subsection*{Your turn}
\begin{itemize}
\raggedright
  \item \emph{Say it:} \textbf{\pa{ਘਰ} \pa{ਵਿੱਚ}}
  \item \emph{Say it:} \textbf{\pa{ਸਕੂਲ} \pa{ਵਿੱਚ}}
  \item \emph{Answer:} after \textbf{\pa{ਕਿੱਥੇ}?}, choose one familiar place chunk
  \item \emph{Check:} did the place come before \textbf{\pa{ਵਿੱਚ}}?
\end{itemize}

\begin{tcolorbox}[breakable,colback=teal!4,colframe=teal!35!black,title={Before you move on}]
How do you say \emph{at home}? (\textbf{\pa{ਘਰ} \pa{ਵਿੱਚ}}.) \emph{At school}? (\textbf{\pa{ਸਕੂਲ} \pa{ਵਿੱਚ}}.) Which comes first, the place or \textbf{\pa{ਵਿੱਚ}}? (\textbf{The place.})

Sources: \href{https://en.wiktionary.org/wiki/\%E0\%A8\%B5\%E0\%A8\%BF\%E0\%A9\%B1\%E0\%A8\%9A}{Wiktionary: \pa{ਵਿੱਚ}}; CDIAL 12042.
\end{tcolorbox}
\section[\texorpdfstring{\pa{ਕਿਤਾਬ} \pa{ਘਰ} \pa{ਵਿੱਚ} \pa{ਹੈ} (kitab ghar vich hai)}{ਕਿਤਾਬ ਘਰ ਵਿੱਚ ਹੈ (kitab ghar vich hai)}]{A thing, a place, then \pa{ਹੈ}}

\label{lesson:PA-C176-location-frame}



\begin{quote}
Say the two pointing pairs: \textbf{\pa{ਇਹ} / \pa{ਉਹ}} and \textbf{\pa{ਇੱਥੇ} / \pa{ਉੱਥੇ}}. Ask \textbf{\pa{ਕਿੱਥੇ}?}. Then say \textbf{\pa{ਘਰ} \pa{ਵਿੱਚ}} and leave \textbf{\pa{ਹੈ}} waiting; Punjabi keeps its verb for the end.
\end{quote}

\begin{grammarlens}[title={Grammar Lens: a thing, a place, then \pa{ਹੈ}}]
\begin{quote}
\textbf{\pa{ਕਿਤਾਬ} \pa{ਘਰ} \pa{ਵਿੱਚ} \pa{ਹੈ।}}
\end{quote}

\begin{quote}
\emph{kitāb ghar vich hai.}
\end{quote}

\begin{quote}
The book is at home.
\end{quote}

The frame has three pieces:

\begin{enumerate}
\raggedright
  \item the thing that is present — \textbf{\pa{ਕਿਤਾਬ}};
  \item its place — \textbf{\pa{ਘਰ} \pa{ਵਿੱਚ}}; and
  \item \textbf{\pa{ਹੈ}} at the end.
\end{enumerate}

Punjabi does not need a separate word matching English \emph{there} in \emph{there is}. The placed thing plus \textbf{\pa{ਹੈ}} states that the thing is present in that location.
\end{grammarlens}

\begin{grammarlens}[title={Grammar Lens: pointing without a postposition}]
When the location itself is \textbf{here} or \textbf{there}, the familiar place word can stand directly before \textbf{\pa{ਹੈ}}:

\begin{quote}
\textbf{\pa{ਦੁਕਾਨ} \pa{ਇੱਥੇ} \pa{ਹੈ।}} — \emph{dukān itthe hai.} — The shop is here.
\end{quote}

\begin{quote}
\textbf{\pa{ਦੁਕਾਨ} \pa{ਉੱਥੇ} \pa{ਹੈ।}} — \emph{dukān utthe hai.} — The shop is there.
\end{quote}

The second sentence is the short practical answer this chapter is building: \textbf{\pa{ਦੁਕਾਨ} \pa{ਉੱਥੇ} \pa{ਹੈ}}.
\end{grammarlens}

\subsection*{Your turn}
\begin{itemize}
\raggedright
  \item \emph{Ask:} \textbf{\pa{ਕਿਤਾਬ} \pa{ਕਿੱਥੇ} \pa{ਹੈ}?}
  \item \emph{Answer:} \textbf{\pa{ਕਿਤਾਬ} \pa{ਘਰ} \pa{ਵਿੱਚ} \pa{ਹੈ।}}
  \item \emph{Ask:} \textbf{\pa{ਦੁਕਾਨ} \pa{ਕਿੱਥੇ} \pa{ਹੈ}?}
  \item \emph{Answer:} \textbf{\pa{ਦੁਕਾਨ} \pa{ਉੱਥੇ} \pa{ਹੈ।}}
  \item \emph{Switch:} point near and answer \textbf{\pa{ਦੁਕਾਨ} \pa{ਇੱਥੇ} \pa{ਹੈ।}}
\end{itemize}

\begin{tcolorbox}[breakable,colback=teal!4,colframe=teal!35!black,title={Before you move on}]
Ask where the book is and answer that it is at home. Then ask where the shop is and answer that it is there. Keep \textbf{\pa{ਹੈ}} last.

This is a small curriculum location exchange against a partial inventory. It is not an A1 mock, not A1 readiness, and not evidence of independent human validation.
\end{tcolorbox}
